\section{WC 28/7/25}
\subsection*{Overview:}
This week I focused on Building my Product video. 

\subsection*{Achievements, Thoughts and Ideas:}
To begin I started by making a 3d Model of a bedroom in blender (a 3d modelling / animating package I have prior experience with) to act as the location for my first scene. I began by blocking out the rough structure I would need and then continued to add details. Whilst working on this model I began to see myself spending a lot of time on small details (such as modelling curtains to direct the light) which wouldn't result in a significant improvement to the final video. To maximise my time spent I decided to stop building the room model and instead Learnt about a new to me technology called lidar room scanning in which I could use my phone to scan my own room and then import that scan for use within my 3d blender scene. 

At this point i switched from my bedroom to my living room as i felt it had more detail and a structure that would be more interesting to animate. To take the scan I first switched off all the lights in the room (as any artificial light would remove my ability to alter the lighting within my 3d scene), guided my dog out of the room (this was a perhaps the most challenging aspect of my EPQ and required many treats and cuddles) and then slowly moved my phone around the room capturing every small detail. 

% \begin{figure}[h]
%     \centering
%     \includegraphics[width=0.8\textwidth]{images/cleaned_scan.png}
%     \caption{capturing the scan of my living room}
%     \label{fig:cleaned_scan}
% \end{figure}

Once I had captured the scan I imported it into blender and began a process called clean-up. As no 3D scan is perfect I had to remove some jagged edges, some small details such as bunting (which was not successfully captured) and replace some elements such as windows and mirrors. These two elements capture / reflect light in reaction to what happens in the 3d scenes so if i for example add a table in front of a mirror I want the mirror to reflect that table. Additionally, I am going to replace the sky with an artificial light source that will need to be seen through the windows.

% \begin{figure}[h]
%     \centering
%     \includegraphics[width=0.8\textwidth]{images/cleaned_scan.png}
%     \caption{The cleaned scan of my living room}
%     \label{fig:cleaned_scan}
% \end{figure}

I then began to add lighting to the scene, I used a combination of point lights (placed as to simulate the rooms main lights + lamps) and a HDRI (High Dynamic Range Image - a 3d image of a sky that contains both data to light objects in my 3d scene and an image that can be seen through the window). I also added a sun lamp to add some additional directional light for atmospheric effect. 
